\par\bigskip\Needspace{9\baselineskip}
\pdfbookmark[2]{Zweite Methode, [3]}{ceres-second-3}
\begin{center}[3.]\end{center}
\begingroup\small
\noindent\hspace*{\parindent}[Der Ausdruck unter dem Integralzeichen ist eine periodische Function der beiden Argumente {\AstroSymbols ⚳} und {\AstroSymbols ♃}, resp. der beiden mittlern Anomalien \(M\) und \(M'\) und kann in eine Reihe nach Vielfachen derselben entwickelt werden. Diese Entwickelung wurde einzeln für jedes Glied nach der in der Theoria interpolationis methodo nova tractata vorbereiteten Methode\refstepcounter{footnote}\hyperlink{ceresIII3-note}{\textsuperscript{*)}}\begingroup
\makeatletter\long\def\@makefntext#1{\noindent[*) #1}\makeatother
\footnotetext{\hypertarget{ceresIII3-note}{}Band III, Seite 265; vgl. auch die Entwickelungen bei den Störungen der Pallas.]}
\endgroup{} vorgenommen und soll hier als Beispiel für das Glied
\[
3)\quad \mathrm{const.}\,rr'Q''\sin(3D+\text{\AstroSymbols ⚳}+3\epsilon'-2\alpha-\AscendingNode),\qquad
\text{wo }\mathrm{const.}=\frac{\mu a\sin\mathrm{J}}{4\BareRoot{(1-ee)}.\sin i}\mathpunct{\raisebox{.45ex}{,}}
\]
\[
\text{numerisch}\qquad\log\mathrm{const.}=2{,}073\ 9541\text{ in Secunden,}
\]
mitgetheilt werden. Zwecks dieser Entwickelung wurde der Umkreis in 10 Theile getheilt und die numerischen Werthe der zu entwickelnden Functionen für die Werthe: \(M\) resp. \(M'=\pm18^{0},\ \pm54^{0},\ \pm90^{0},\ \pm126^{0},\ \pm162^{0}\) berechnet, wo \(M\) und \(M'\) von der Sonnenferne gezählt sind. Setzt man für den Augenblick, um die weiter unten Art. 2 der allgemeinen Störungen der Pallas, erste Rechnung, gegebenen Formeln benutzen zu können, \(m=90^{0}-M\), und will man eine Function \(T\) in die Reihe:
\SourceLeaf{410}
\[
\begin{aligned}
T={}&\begin{aligned}[t]&\alpha_0+\alpha_1\cos m+\alpha_2\cos2m+\ldots\\&+\beta_1\sin m+\beta_2\sin2m+\ldots\end{aligned}\\[4pt]
={}&\begin{aligned}[t]&\alpha_0-\beta_1\cos M-\alpha_2\cos2M+\ldots\\&+\alpha_1\sin M-\beta_2\sin2M-\ldots\end{aligned}
\end{aligned}
\]
entwickeln, und sei
\[
\begin{aligned}
\text{für die Werthe }m={}&0,\ 36^{0},\ 72^{0},\ \ldots\ 288^{0},\ 324^{0}\\
T={}&A_0,\ A_1,\ A_2,\ \ldots\ A_8,\ A_9
\end{aligned}
\]
gegeben, so ist nach den a. a. O. entwickelten Formeln
\[
10\alpha_0=A_0+A_2+A_4+A_6+A_8-(A_1+A_3+A_5+A_7+A_9)
\]
\[
\left.\begin{aligned}
5\alpha_n={}&A_0-A_5+(A_1+A_9-A_4-A_6)\cos na\\
&+(A_2+A_8-A_3-A_7)\cos2na\\[3pt]
5\beta_n={}&(A_1-A_9+A_4-A_6)\sin na\\
&+(A_2-A_8+A_3-A_7)\sin2na
\end{aligned}\right\}\quad n\text{ ungerade}
\]
\[
\left.\begin{aligned}
5\alpha_n={}&A_0+A_5+(A_1+A_9+A_4+A_6)\cos na\\
&+(A_2+A_8+A_3+A_7)\cos2na\\[3pt]
5\beta_n={}&(A_1-A_9-A_4+A_6)\sin na\\
&+(A_2-A_8-A_3+A_7)\sin2na
\end{aligned}\right\}\quad n\text{ gerade}
\]
\[
a=36^{0}.
\]

Das zu entwickelnde Glied 3) wurde nun wie folgt zerlegt:
\[
\begin{aligned}
&rr'Q''\sin(3D+\text{\AstroSymbols ⚳}+3\epsilon'-2\alpha-\AscendingNode)=\\
&\quad r'\cos3\epsilon'\left\{\begin{aligned}&r\cos2\alpha.\,Q''\sin(3D+\text{\AstroSymbols ⚳}-\AscendingNode)\\&-r\sin2\alpha.\,Q''\cos(3D+\text{\AstroSymbols ⚳}-\AscendingNode)\end{aligned}\right\}\\[4pt]
&\quad+r'\sin3\epsilon'\left\{\begin{aligned}&r\cos2\alpha.\,Q''\cos(3D+\text{\AstroSymbols ⚳}-\AscendingNode)\\&+r\sin2\alpha.\,Q''\sin(3D+\text{\AstroSymbols ⚳}-\AscendingNode)\end{aligned}\right\}
\end{aligned}
\]
und es wurden zunächst die Ausdrücke \(\displaystyle r\genfrac{}{}{0pt}{}{\cos}{\sin}2\alpha.\,Q''\) für alle Werthepaare
\[
M\text{ resp. }M'=\pm18^{0},\ \pm54^{0},\ \pm90^{0},\ \pm126^{0},\ \pm162^{0}
\]
berechnet.]
\par\endgroup
\begin{center}\rule{.30\linewidth}{.4pt}\end{center}
